\documentclass[12pt]{article}
\usepackage{amsmath, amssymb, geometry, enumitem}
\geometry{margin=1in}
\setlength{\parindent}{0pt}
\setlength{\parskip}{1em}

\title{ECON 101: Macroeconomics \vspace{5pt}\\ Quiz 6}
\author{Instructor: Muhebullah Karimzada}
\date{December 03,  2025}

\begin{document}
\maketitle

\section*{Multiple Choice Questions (1 mark each)}
\textit{Each question is worth 1 mark. Choose the best answer.}

\setlist[enumerate]{itemsep=2pt, topsep=2pt}

\textbf{1.} The Bank of Canada's primary monetary policy instrument is:
\begin{enumerate}[label=(\Alph*)]
    \item Open-market purchases
    \item The required reserve ratio
    \item The target overnight interest rate
    \item Foreign exchange interventions
\end{enumerate}

\textbf{2.} If the Bank of Canada lowers the target overnight rate, the most immediate effect is:
\begin{enumerate}[label=(\Alph*)]
    \item Decrease in bank reserves
    \item Increase in short-term market interest rates
    \item Decrease in consumption spending
    \item Increase in settlement balances and lower short-term rates
\end{enumerate}

\textbf{3.} An open-market \emph{sale} of government securities by the Bank of Canada will:
\begin{enumerate}[label=(\Alph*)]
    \item Increase the money supply and reduce interest rates
    \item Reduce bank reserves and raise interest rates
    \item Increase bond prices and lower yields
    \item Shift aggregate demand to the right
\end{enumerate}

\textbf{4.} The monetary policy transmission mechanism begins with:
\begin{enumerate}[label=(\Alph*)]
    \item Changes in consumption and investment
    \item Movements in the exchange rate
    \item A change in the policy interest rate
    \item Shifts in long-run aggregate supply
\end{enumerate}

\textbf{5.} If the Bank of Canada raises interest rates, the Canadian dollar tends to:
\begin{enumerate}[label=(\Alph*)]
    \item Depreciate, increasing exports
    \item Appreciate, reducing net exports
    \item Stay constant due to inflation targeting
    \item Fall against the US dollar but rise against the euro
\end{enumerate}

\textbf{6.} Which of the following is \emph{part} of the Bank of Canada's inflation-control strategy?
\begin{enumerate}[label=(\Alph*)]
    \item A fixed money supply growth rate
    \item A 2\% inflation target midpoint
    \item Exchange-rate targeting
    \item Wage growth ceilings
\end{enumerate}

\textbf{7.} Contractionary monetary policy in the AD--AS model will:
\begin{enumerate}[label=(\Alph*)]
    \item Shift AD to the right and raise inflation
    \item Shift SRAS to the left and raise inflation
    \item Shift AD to the left and lower real GDP in the short run
    \item Shift LRAS to the right and increase long-run growth
\end{enumerate}


\newpage

\section*{Problem (13 marks)}



The Bank of Canada considers the following reaction function when setting its target overnight interest rate:

\[
i = 1 + 1.2(\pi - 2) + 0.4(Y - Y^*),
\]

where:
\begin{itemize}
    \item \( i \) = target overnight rate (percent),
    \item \( \pi \) = actual inflation rate (percent),
    \item \( 2 \) = inflation target midpoint (percent),
    \item \( Y - Y^* \) = output gap (percent, positive if economy is above potential).
\end{itemize}

New data show:
\[
\pi = 3.1\%, \qquad Y - Y^* = +0.8\%.
\]



\begin{enumerate}
    \item[(a)] Compute the recommended Bank of Canada target overnight interest rate using the reaction function. Show all steps. (4 marks)
    \item[(b)] Explain whether the resulting policy is expansionary or contractionary. (2 marks)
    \item[(c)] Describe, step by step, how an increase in the policy rate affects consumption, investment, exchange rate, and net exports. (4 marks)
    \item[(d)] Using the AD--AS model, explain how this policy helps bring inflation back toward the 2\% target. (3 marks)
\end{enumerate}

\hrule
\vspace{1em}

\begin{center}
\textit{End of Quiz}
\end{center}

\end{document}
